% Determinac ión Del Perfil Del Autor En Contextos De Clasificación Anticipada De Textos Para múltiples Dominios =>

% 1. Determinación Del Perfil Del Autor
% 2. Redes sociales, Clasificación Anticipada
% 3. múltiples Dominios (de riesgo, depresión, anorexia, self-harm)

% Motivación [plantear el gap: el mismo que el primer paper, "no existe" un modelo que cumpla con esos 3 requisitos de forma natural e integrada]

% (estructurar como la tiene Mike directamente, más fácil + Publicaciones Resultantes de este Trabajo)

% Este capítulo sirve como una introducción general donde se realiza la presentación del problema abordado, 
% se consideran los principales antecedentes relacionados con el tema que han sido encontrados en la bibliografía consultada, 
% se enuncian los objetivos perseguidos y finalmente se presenta la organización del  informe.

Las redes sociales son, en esencia, una plataforma para facilitar que las personas expresen y compartan sus ideas, pensamientos, y opiniones.
Son, además, un lugar de encuentro que permite que las personas se conecten entre sí, como lo hemos hecho durante miles y miles de años, pero eliminando las limitaciones espaciales, y temporales, inherentes a los métodos tradicionales de comunicación.

Durante los últimos años, las redes sociales vienen disfrutando de una gran popularidad que aumenta año a año. Por ejemplo, Facebook, la red social más grande del mundo, posee más de 2.700 millones de usuarios activos mensuales\cite{facebook2020}; Twitter alberga alrededor de 275 millones de usuarios activos mensuales que publican un promedio de 500 millones de tweets por día \cite{twitter2020}; Reddit, superando a Twitter, posee más de 430 millones de usuarios activos mensuales\cite{reddit2020}; además, para el año 2012 ya se estimaba que había más de 181 millones de blogs activos en todo el mundo\cite{nielsen2012}.

En otras palabras, a cada instante las personas están generando contenido propio, en enormes cantidades, expresando y compartiendo sus ideas, pensamientos y opiniones. Sin embargo, más allá de algunos datos puntuales, como ser su nombre de usuario, la mayoría de las veces no se tiene acceso a más información sobre el autor de este contenido.
No obstante, numerosas son las razones por las que muchas veces es fundamental conocer cierta información relevante de los usuarios de estas redes sociales.
Por ejemplo, desde el punto de vista del marketing existe un interés en conocer la identidad y características demográficas de los distintos usuarios, con la intención de orientar la publicidad para que su efectividad sea la mayor posible\cite{bentolila2011system};
en el área de interacción humano-computadora, es fundamental conocer las características específicas de cada persona para poder mostrar una interfaz acorde a la personalidad de cada individuo\cite{andres2015towards};
desde la lingüística forense\cite{aronoff2017handbook}, también es importante hacer uso del conocimiento lingüístico para analizar el contenido de los usuarios con el fin de detectar y prevenir  actos ilícitos o engañosos, como ser extorsiones o, incluso, acosos sexuales\cite{hall2007profile, escalante2015early}.

Debido al volumen abrumador de contenido que minuto a minuto generan millones de usuarios,  realizar este tipo de análisis de forma manual sería completamente impensable. Por este motivo, ha surgido una creciente necesidad de realizar este tipo de análisis utilizando modelos y técnicas científico-computacionales.
Por ejemplo, dentro de las ciencias computacionales, más precisamente dentro del campo denominado como \emph{\ac{PLN}}, se conoce como \emph{análisis de autoría} al estudio de diversas cuestiones vinculadas al autor de un contenido textual\cite{indurkhya2010handbook}.
Más precisamente, el \emph{análisis de autoría} estudia e investiga procesos y técnicas computacionales para examinar las características de un texto con el fin de obtener información y/o conclusiones sobre su autor\cite{el2014authorship}.
A su vez, en la literatura\cite{zheng2003authorship,abbasi2005applying,zheng2006framework}, típicamente se divide el \emph{análisis de autoría} en dos áreas principales de investigación, a saber, \emph{\ac{AA}} y \emph{\ac{DPA}}\footnote{``Author Profiling'' en inglés}.
Si bien ambos parten del análisis del texto, cada área busca responder a distintas preguntas sobre su autor.
Así, la \ac{AA} consiste en determinar, entre un conjunto de posibles autores, quién es efectivamente el autor del mismo\cite{stamatatos2009survey}, mientras que la \ac{DPA} busca determinar diversas características sobre el autor\cite{argamon2009automatically}.
La hipótesis clave detrás de la \ac{DPA} es que la forma en que escribimos revela nuestro comportamiento y nuestra personalidad. 
Un ejemplo clásico de tarea de \ac{DPA} sería la determinación del perfil sociodemográfico del autor en base al contenido que éste genera, como ser su idioma nativo, género, etnia, edad, ubicación, ocupación, o su nivel socio-económico\cite{koppel2005determining,schler2006effects, corney2002gender}.
Sin embargo, también se han realizado esfuerzos para determinar perfiles centrados en otro tipo de aspectos no tan directos, como ser el nivel de bienestar\cite{schwartz2013characterizing}, la ideología política\cite{koppel2009automatically} o, incluso, rasgos de personalidad como la extraversión o el neuroticismo\cite{argamon2005lexical,mairesse2006words,rangel2015overview}.

\section{Planteamiento del Problema}

El lenguaje es un poderoso indicador no sólo de la personalidad y el estado social o emocional de un individuo, sino también de su salud mental. Por ejemplo, desde la Psicología, el vínculo entre el uso del lenguaje y los trastornos clínicos se ha estudiado durante décadas\cite{pennebaker2003psychological}.
En este contexto, una tarea de \ac{DPA} de especial interés, debido al enorme impacto que puede tener en la vida de la persona, es la determinación del perfil psicológico del autor, vinculado a posibles trastornos mentales.
En particular, una problemática que recientemente está ganando mucha atención es la \ac{DAR} en las redes sociales, vinculados a diversos trastornos mentales, como son la depresión, la anorexia o las conductas autodestructivas\cite{losada2016test,losada2017erisk,losada2018overview,losadaetal2019}. %, ya que puede significar la diferencia entre vivir o morir.
Más en detalle, esta tarea de \ac{DAR} consiste en procesar secuencialmente el contenido que cada usuario va generando hasta determinar en algún momento, \emph{lo antes posible},\footnote{De allí  ``anticipada''.} si el usuario está en riesgo debido al padecimiento de algún tipo de trastorno. El objetivo central que impulsa el reconocimiento temprano de estas situaciones de riesgo es garantizar que la persona reciba la atención y el apoyo social que necesita a tiempo, ya que de no ser así, podría incluso hasta costarle la vida.
Sin embargo, abordar esta tarea es extremadamente desafiante ya que, a diferencia de las tareas de \ac{DPA} convencionales, aquí los sistemas no sólo deben determinar el perfil (psicológico) del autor al clasificar su posible trastorno, sino que, además, deben ser capaces de decidir si la información procesada hasta el momento es suficiente para hacerlo con ciertas garantías o, si de lo contrario, se debe esperar a que el usuario genere más contenido, con el riesgo que esto implique.
Más aún, al igual que con cualquier otra aplicación crítica en salud, finanzas o seguridad nacional, este es un dominio que se beneficiaría enormemente con sistemas que no sólo hagan predicciones correctas, sino que también faciliten la comprensión de cómo las llevaron a cabo.
De hecho, desarrollar estos sistemas con técnicas y modelos de aprendizaje automático tradicionales plantea aspectos realmente desafiantes para la Inteligencia Artificial.
En síntesis, podemos identificar al menos tres aspectos que idealmente estos modelos deberían ser capaces de soportar: \emph{capacidad de trabajar incrementalmente}, \emph{soportar algún mecanismo de toma de decisión anticipada} y \emph{capacidad de justificar sus decisiones}.

Sin embargo, como se resaltará y analizará con más detalle en este trabajo, la mayoría de los modelos de aprendizaje automático supervisados resultan no ser adecuados para abordar este tipo de escenarios, ya sea porque o bien funcionan como cajas negras o no son capaces de trabajar incrementalmente. De hecho, y hasta donde sabemos, ningún modelo de clasificación de texto es capaz de soportar los tres aspectos mencionados anteriormente, de una manera natural e integrada.


\section{Preguntas de Investigación}

%[el problema de arriba deben ser los 3 puntos necesarios: incremental, early y interpretabilidad => entonces acá debo preguntar: ¿los clasificadores actuales son aptos? ¿es posible diseñar un nuevo clasificador que cumpla con estos tres aspectos? + preguntas propias de los experimentos]

A lo largo del presente trabajo, intentamos responder a las siguientes preguntas de investigación:

\begin{itemize}
    \item[\textbf{P1:}] \emph{¿Qué características deberían idealmente tener los clasificadores de texto para poder abordar, eficiente e integralmente, tareas de detección anticipada de riesgos en entornos dinámicos, como el de las redes sociales?}
    \item[\textbf{P2:}] \emph{¿Cómo podemos diseñar un nuevo clasificador de texto que naturalmente incorpore dichas características de una forma unificada?}
    \item[\textbf{P3:}] \emph{¿Qué tan efectivo y robusto podría llegar a ser este modelo en determinar anticipadamente el perfil psicológico del autor para distintos dominios, como son anorexia, depresión o conductas autodestructivas?}
\end{itemize}

\section{Objetivos}
\label{sec:objetivos}

\subsection{Objetivo General}

Diseñar, desarrollar y evaluar un nuevo modelo de aprendizaje automático para clasificación de texto que sea naturalmente apto para abordar tareas de detección anticipada de riesgos en entornos dinámicos, como el de las redes sociales, de forma eficiente e integral.

\subsection{Objetivos Específicos}

\begin{itemize}
    \item[\textbf{O1:}] Analizar y detectar las características que  un clasificador de texto debería poseer para poder abordar tareas de detección anticipada de riesgo en ambientes dinámicos y masivos como lo son las redes sociales.
    \item[\textbf{O2:}] Diseñar e implementar un nuevo clasificador de texto que incorpore dichas características de forma natural.
    \item[\textbf{O3:}] Evaluar el desempeño y la robustez del modelo propuesto en tareas de detección anticipada de depresión, anorexia y  conductas autodestructivas.
    % \item[\textbf{O4:}] Analizar, proponer, desarrollar y evaluar al menos una posible mejora del modelo original.
    \item[\textbf{O4:}] Crear un paquete de código libre con la implementación del modelo propuesto y ponerlo a disposición de la comunidad para que pueda ser utilizado libremente.
\end{itemize}

% \section{Justificación de la Solución}
\section{Descripción del Trabajo Propuesto y su Relevancia Científica}

En este trabajo se propone, en concordancia con el objetivo planteado en la sección anterior, un nuevo modelo de aprendizaje automático al cual hemos denominado como SS3. A diferencia de los modelos de clasificación convencionales, SS3 no procesa la entrada completa como un único vector indivisible de atributos, $\pmb{x}=<x_1,x_2,\dots,x_m>$, sino que fue diseñado para procesarla secuencialmente, elemento a elemento.
Esta característica le permite al modelo poder trabajar incrementalmente, dotándole de la capacidad de poder clasificar la entrada, potencialmente, en cualquier instante de tiempo.
En este sentido, uno podría notar que es una característica compartida con otros modelos que también poseen la capacidad de procesar la entrada de esta manera, como son las Redes Neuronales Recurrentes o el clasificador \emph{Multinomial Naïve Bayes}.
Sin embargo, a diferencia de estos últimos, SS3 posee la capacidad de explicar intuitivamente las razones que llevaron a clasificar la entrada como tal.
Esto se debe a que nuestro modelo fue diseñado para aprender, explícitamente, qué tan valioso es cada uno de los elementos de entrada para dicha clasificación.
Más precisamente, SS3 aprende a asignarle un valor a cada elemento de entrada, $w$, en relación a cada categoría, $c$, utilizando una función especial.
Esta función captura el \emph{valor de confianza} con el que se cree que un elemento $w$ ``es importante'' para una categoría $c$.
Por ejemplo, si suponemos que los elementos individuales de entrada son palabras y que, además, las categorías de interés son $C= \{$\emph{tecnología, salud, alimentos}$\}$, entonces, luego de entrenarlo, SS3 aprendería a valorar las palabras ``sushi'' y ``el'' de una manera similar a lo mostrado a continuación:

\vspace{5mm}
\begin{tabular}{l l@{\hskip 1in}l l}
% \begin{tabular}{@{}l@{}l}
$valor(\text{\emph{``sushi''}}, \text{\emph{tecnología}}) $ & $= 0$ & $valor(\text{\emph{``el''}}, \text{\emph{tecnología}}) $ & $= 0$\\
$valor(\text{\emph{``sushi''}}, \text{\emph{salud}}) $ & $= 0.55$ & $valor(\text{\emph{``el''}}, \text{\emph{salud}}) $ & $= 0$\\
$valor(\text{\emph{``sushi''}}, \text{\emph{alimentos}}) $ & $= 0.75$  & $valor(\text{\emph{``el''}}, \text{\emph{alimentos}}) $ & $= 0$\\
\end{tabular}
\vspace{5mm}

Luego, en su versión más simple, la clasificación de la entrada completa se lleva a cabo seleccionando la categoría que tenga el mayor \emph{valor de confianza} acumulado, sumando el valor de todos los elementos leídos hasta el momento. Más formalmente, siendo $C$ el conjunto de todas las categorías de interés, siendo $\oplus$ alguna operación asociativa,\footnote{Como ser la suma.} y siendo $w_1,w_2,\dots,w_n$ una (secuencia de) entrada de $n$ elementos,\footnote{En la práctica, normalmente $n$ es virtualmente infinito.} SS3 asignará a la entrada la categoría $\hat{c}$ tal que:


$$\hat{c}=\argmax_{c\in C} \bigoplus_{i=1}^{n} valor(w_i,c)$$

\

Para ilustrar esta ecuación con un ejemplo, supongamos que se desea clasificar la pequeña oración ``el sushi sabroso'', y que la operación $\oplus$ elegida es la suma, entonces primero se calcularía la suma de los valores de cada palabra para cada categoría, obteniendo así el valor de confianza acumulada para cada una de ellas, del siguiente modo: 

\begin{tabular}{l@{}l}%p{0.5\linewidth}}
$\text{\emph{confianza para tecnología}}$ & $=  valor(\text{\emph{``el''}}, \text{\emph{tecnología}}) + valor(\text{\emph{``sushi''}}, \text{\emph{tecnología}})$ \\
& \ \ \ $ +\ valor(\text{\emph{``sabroso''}}, \text{\emph{tecnología}}) = 0 + 0 + 0 $\\
& $= 0$\\
$\text{\emph{confianza para salud}}$ & $=  valor(\text{\emph{``el''}}, \text{\emph{salud}}) + valor(\text{\emph{``sushi''}}, \text{\emph{salud}})$ \\
& \ \ \ $ +\ valor(\text{\emph{``sabroso''}}, \text{\emph{salud}}) = 0 + 0.55 + 0.10 $\\
& $= 0.65$\\
$\text{\emph{confianza para alimentos}}$ & $=  valor(\text{\emph{``el''}}, \text{\emph{alimentos}}) + valor(\text{\emph{``sushi''}}, \text{\emph{alimentos}})$ \\
& \ \ \ $ +\ valor(\text{\emph{``sabroso''}}, \text{\emph{alimentos}}) = 0 + 0.75 + 0.80$\\
& $= 1.55$ \\
\end{tabular}

\

Luego, SS3 clasificaría la oración con la categoría que obtuvo el mayor valor de confianza acumulada, es decir ``alimentos'', cuyo valor fue de $1.55$.

\


Notar que esta función $valor(w,c)$ fue especialmente diseñada para valorar cada elemento intentando imitar la manera en la que una persona respondería la pregunta ``¿cuánto vale $w$ para la categoría $c$?'' ---en nuestro ejemplo, ``¿cuánto vale la palabra \emph{sushi} para alimentos?'', ``¿cuánto la palabra \emph{el} para alimentos?'',\footnote{Notar que el modelo aprendió que la palabra ``el'', al ser un pronombre, no tiene valor para ninguna categoría en particular. En su lugar, por ejemplo, Multinomial Naïve Bayes respondería esta pregunta asignando los valores más altos, por ejemplo, a todos los pronombres; ya que a diferencia de SS3, MNB valora a las palabras solamente según su probabilidad de ocurrencia, local, en cada categoría.} y así.
Es esta característica de diseño que dota al modelo, a diferencia de los otros modelos existentes, la capacidad de poder explicar natural e intuitivamente\footnote{Ya que el mismo diseño parte de imitar la forma en la que una persona valoraría cada elemento.} las razones que lo llevaron a clasificar la entrada como tal. En la práctica, estas explicaciones pueden ser visuales, coloreando a cada elemento de entrada con una intensidad proporcional a su valor.\footnote{Esto puede apreciarse en las demos que hemos desarrollado para que los lectores interesados puedan probar SS3 en tiempo real, accediendo a  \url{http://tworld.io/ss3/}.}

\

Finalmente, antes de terminar esta sección, creemos conveniente resaltar la relevancia científica del trabajo propuesto, el cual no se restringe a la presentación de un modelo con el que se podrían abordar tareas de detección anticipada de riesgo con cierto grado de eficiencia, sino que también introduce un nuevo modelo de aprendizaje automático supervisado de propósito general que, hasta donde sabemos, sería el primero en integrar todas las características mencionadas anteriormente.
Dada la naturaleza de ``caja blanca'' de SS3, este clasificador podría permitirle a los investigadores y profesionales implementar modelos de clasificación de texto más transparentes y, por lo tanto, más confiables.\footnote{Esto podría ser especialmente útil, por ejemplo, para quienes trabajen con problemas sensibles que involucren a personas.} Lo que, creemos, potencialmente podría llegar a ser un aporte significativo al campo del aprendizaje automático en su forma más general.
% Finalmente, antes de terminar esta sección, creemos conveniente resaltar la relevancia científica del trabajo propuesto, que no se limita a la presentación de un modelo con el que se podrían construir sistemas de detección anticipada de riesgo simples y eficientes, como lo sugieren los resultados experimentales, sino que también introduce un nuevo modelo de aprendizaje automático supervisado de propósito general que, hasta donde sabemos, es el primero en integrar todas las características mencionadas anteriormente.
% Dada la naturaleza de ``caja blanca'' de SS3, este clasificador podría permitirle a los investigadores y profesionales implementar modelos de clasificación de texto más transparentes y, por lo tanto, más confiables.\footnote{Esto podría ser especialmente útil, por ejemplo, para quienes trabajen con problemas sensibles que involucren a personas.} Lo que, creemos, potencialmente puede ser un aporte  significativo al campo del aprendizaje automático como tal.

\section{Publicaciones Resultantes de este Trabajo}

Como consecuencia de esta investigación doctoral se han publicado los siguientes artículos:

\

\noindent
\textbf{Revistas Internacionales (JCR)}

\begin{itemize}
    \item Burdisso, S. G.; Errecalde, M.; Montes-y-Gómez, M. (2019). A text classification framework for simple and effective early depression detection over social media streams. \emph{Expert Systems with Applications}, vol. 133, pp. 182-197. %DOI:  \url{10.1016/j.eswa.2019.05.023}
    \begin{itemize}
        \item \textbf{JCR Impact Factor:} 5.452
        \item \textbf{SJR Quartiles:} Artificial Intelligence (Q1); Computer Science Applications (Q1); Engineering (Q1)
    \end{itemize}
    \item Burdisso, S. G.; Errecalde, M.; Montes-y-Gómez, M. (2020). $\tau$-SS3: a text classifier with dynamic n-grams for early risk detection over text streams. \emph{Pattern Recognition Letters}, vol. 138, pp. 130-137. %DOI: \url{10.1016/j.patrec.2020.07.001}
    \begin{itemize}
        \item \textbf{JCR Impact Factor:} 3.255
        \item \textbf{SJR Quartiles:} Computer Vision and Pattern Recognition (Q1); Software (Q1); Artificial Intelligence (Q2); Signal Processing (Q2)
    \end{itemize}
    % \item Burdisso, S. G.; Errecalde, M.; Montes-y-Gómez, M. (2020). PySS3: A Python package implementing a new model for text classification with visualization tools for Explainable AI. \emph{Computing in Science \& Engineering} (en revisión).
    % \begin{itemize}
    %     \item \textbf{JCR Impact Factor:} 1.935
    %     \item \textbf{SJR Quartiles:} Engineering (Q1); Computer Science (Q2);
    % \end{itemize}
\end{itemize}

\noindent
\textbf{Revistas Nacionales}

\begin{itemize}
    \item Burdisso, S. G.; Errecalde, M.; Montes-y-Gómez, M. (2020). Using Text Classification to Estimate the Depression Level of Reddit Users. \emph{Journal of Computer Science \& Technology}, vol. 21, no. 1, pp. 1–10. %DOI:  \url{10.1016/j.eswa.2019.05.023}
\end{itemize}

\noindent
\textbf{Capítulos de Libro}

\begin{itemize}
    \item Burdisso, S. G.; Cagnina L. C.; Errecalde, M.; Montes-y-Gómez, M. (a ser publicado a mediado del 2021). Two Simple and Domain-independent Approaches for Early Detection of Anorexia. \emph{Early Risk Detection of Psychological Conditions: the eRisk Experience}, Springer.
\end{itemize}

\noindent
\textbf{Conferencias Internacionales}

\begin{itemize}
    \item Burdisso, S. G.; Errecalde, M.; Montes-y-Gómez, M. (2019). UNSL at eRisk 2019: a Unified Approach for Anorexia, Self-harm and Depression Detection in Social Media.  \emph{Working Notes of CLEF 2019}, 2380, 2019, Lugano, Suiza.
    \item Funez, D.; Ucelay M. J.; Villegas M. P.; Burdisso, S. G.; Cagnina L. C.;  Montes-y-Gómez, M.; Errecalde, M. (2018). UNSL's participation at eRisk 2018 Lab. \emph{Working Notes of CLEF 2018}, 2125, Avignon, Francia.
\end{itemize}

\noindent
\textbf{Congresos Nacionales}

\begin{itemize}
    \item Burdisso, S. G.; Errecalde, M.; Montes-y-Gómez, M. (2019).
    Towards Measuring the Severity of Depression in Social Media via Text Classification. \emph{Actas del XXV Congreso Argentino de Ciencias de la Computación (CACIC 2019)}, pp. 577-588. 2019. Río Cuarto, Córdoba, Argentina.
    \item Cagnina L. C.; Ferretti E.; Villegas M. P.; Ucelay M. J.; Burdisso, S. G.; Funez, D.; Errecalde, M. (2016). Minería de Textos y de la Web. \emph{XVIII Workshop de Investigadores en Ciencias de la Computación, WICC 2016}, Entre Ríos, Argentina.
\end{itemize}


\section{Estructura del Documento}

Este documento de tesis consta de un total de ocho capítulos agrupados en cuatro partes en las que, respectivamente, se introduce el trabajo, se revisan los antecedentes teóricos, se presentan las contribuciones realizadas y se brindan las conclusiones. Más detalladamente, luego de este capítulo introductorio, el resto del documento se encuentra estructurado de la siguiente manera:

\begin{itemize}
    \item La \autoref{part2} introduce los antecedentes teóricos necesarios para que este trabajo de tesis sea lo más autocontenido  posible. Esta parte se divide en los siguientes capítulos:
    \begin{itemize}
        \item El \autoref{cap2} sirve a modo de introducción al Aprendizaje Automático Supervisado, en el que se introducen conceptos vinculados a la clasificación, se presentan los clasificadores más populares y se describen las medidas de desempeño utilizadas para evaluar el rendimiento de los clasificadores.
        \item El \autoref{cap3} tiene un propósito doble, servir como una breve introducción a la clasificación de texto y, finalmente, introducir al lector la problemática abordada, esta vez definiéndola dentro de un marco teórico más técnico.
    \end{itemize}
    \item La \autoref{part3} presenta las principales contribuciones de este trabajo de tesis, en relación a los objetivos planteados en la \autoref{sec:objetivos}, y se compone de los siguientes capítulos:
    \begin{itemize}
        \item El \autoref{cap:ss3} introduce, en el contexto de detección anticipada de depresión, el nuevo modelo de aprendizaje automático desarrollado en relación al objetivo general de esta tesis doctoral.
        \item El \autoref{cap:tss3} propone una extensión del modelo original y luego la evalúa, junto al original, en tareas de detección anticipada de depresión y anorexia.
        \item El \autoref{cap:eRisk2019} describe nuestra participación, utilizando el modelo propuesto, en el eRisk 2019 del CLEF, específicamente en las tareas de detección anticipada de anorexia y conductas autodestructivas.
        \item El \autoref{cap:pyss3} introduce y describe el paquete Python creado con el fin de que la comunidad pueda utilizar, de una manera sencilla y transparente, el modelo desarrollado.
    \end{itemize}
    \item La \autoref{part4} contiene el \autoref{cap:conclusion}, en el que se resumen las principales conclusiones derivadas de este trabajo y se sugieren posibles lineas de trabajo a futuro.
    \hfill$\square$
\end{itemize}
